% !PRINTED 472
\weberpage{472}

%<<CONT
man $h_x = 0$ setzt. Die Gr\"ossen $u_x, g_x, h_x$ h\"angen nur von den
Variablen $x_1, x_2, x_3$ ab.

Nun l\"asst sich, wie \gesperrt{Klein} a.\,a.\,O. ${}^1$) bemerkt hat, die L\"osung
der allgemeinen Resolvente auf die der speciellen zur\"uckf\"uhren,
wenn man die Wurzel einer \gesperrt{biquadratischen Gleichung}
adjungirt.

Um dies nachzuweisen, f\"uhren wir neben dem Punkte $(x)$
einen zweiten Punkt $(y)$ ein und bilden die Polare von $f$:
\weq{3}{f_1(x,y) = y_1f'(x_1) + y_2f'(x_2) + y_3f'(x_3).}
Wenn wir dann $(x)$ und $(y)$ gleichzeitig derselben Substitu\-tion
der Gruppe $G_{168}$ unterwerfen, so bleibt die Function $f_1(x,y)$
unge\"andert (Bd.\ I, \S.\,60). Wir nehmen nun den Punkt $(x)$ beliebig
 an, verlangen aber von dem Punkte $(y)$, dass er gleichzeitig
 auf der Grundcurve und auf der Polaren des Punktes $(x)$
liegen soll, dass also gleichzeitig
\weq{4}{f(y_1,y_2,y_3) = 0, \quad f_1(x,y) = 0}
sein soll. Dann entsprechen jedem Punkte $x$ vier Punkte $y$, und
wenn wir f\"ur $(x)$ einen mit ihm verbundenen Punkt setzen, so
geht jeder dieser vier Punkte $(y)$ gleichfalls in einen verbundenen
Punkt \"uber. Die Function $h_y$ ist jetzt $=0$ und $g_y$ ist eine vierwerthige
 Function des Punktes $(x)$. Symmetrische Functionen
dieser vier Werthe bleiben unge\"andert durch die Substitutionen
von $G_{168}$, und folglich ist $g_y$ \gesperrt{Wurzel einer biquadratischen
Gleichung, deren Co\"efficienten rational von den $g_x, h_x$
abh\"angen}. Die Wurzel dieser biquadratischen Gleichung muss
adjungirt werden.

Die Function $u_y$ ist Wurzel der speciellen Resolvente
\weq{5}{R(u_y, g_y, 0) = R_y = 0.}
Wir bezeichnen nun die vier zu demselben $x$ geh\"origen Punkte
$y$ mit $y, y', y'', y'''$ und bilden die symmetrische Function dieser
vier Punkte
\weq{6}{(t-u_y)(t-u_{y'})(t-u_{y''})(t-u_{y'''}) = \Phi(t)}
f\"ur ein unbestimmtes $t$. Diese Function bleibt unge\"andert bei
allen Substitutionen der Gruppe $G_{24}$, und ist also rational durch
$u_x, g_x, h_x$ ausdr\"uckbar.

\par\noindent\rule{0.3\linewidth}{0.4pt}\par
${}^1$) Mathematische Annalen, Bd.\ XV, S.\ 280.
